\documentclass[11pt]{article}
%
%
% Retirez le caractere "%" au debut de la ligne ci-dessous si votre
% editeur de texte utilise des caracteres accentues
% \usepackage[latin1]{inputenc}
\setlength{\headheight}{14pt}
%
% Retirez le caractere "%" au debut des lignes ci-dessous si vous
% utiisez les symboles et macros de l'AMS
 \usepackage{amsmath}
\usepackage{amsfonts}
\usepackage{amsthm}
\newtheorem{lemma}{Lemma}
\newtheorem{proposition}{Proposition}

\usepackage{fancyhdr}
\pagestyle{fancy}
\usepackage{graphicx}
\usepackage{color}

%\usepackage{biblatex}
%\addbibresource{Biblio.bib}

\usepackage[left=2.5cm,right=2.5cm,top=2.5cm,bottom=2.5cm]{geometry}
\setlength{\parskip}{6pt}
\usepackage{lmodern}
\usepackage[T1]{fontenc}
\usepackage[utf8]{inputenc}
\usepackage[english]{babel}
%\usepackage{caption}


\usepackage{amsopn}             % pour \DeclareMathOperator
\usepackage{amssymb}            % pour \leqslant

\DeclareMathOperator{\cov}{cov}
\DeclareMathOperator{\cor}{cor}
\DeclareMathOperator{\var}{var}
\DeclareMathOperator{\sign}{sign}
\DeclareMathOperator{\sumss}{SS}
\DeclareMathOperator{\diag}{diag}
\DeclareMathOperator{\rank}{rank}
\DeclareMathOperator*{\argmin}{arg\,min}
\DeclareMathOperator*{\argmax}{arg\,max}
\DeclareMathOperator*{\minimize}{minimize}
\DeclareMathOperator{\trace}{\mathrm{trace}}

\newcommand{\X}{\textbf{X}}
\newcommand{\U}{\textbf{U}}
\newcommand{\V}{\textbf{V}}
\renewcommand{\u}{\textbf{u}}
\renewcommand{\v}{\textbf{v}}
\newcommand{\HH}{\textbf{H}}
\newcommand{\EE}{\textbf{E}}

\begin{document}
\renewcommand{\headrulewidth}{0pt}


%--------------------------------------------------------------------------

\begin{center}
{\Large
	{\sc Optimal Canonical factorization of multiblock data}
}
\bigskip

 \underline{Jean-Michel Galharret}$^{1}$ \& Stephanie Bougeard $^{2}$ \& Katrijn Van Deun$^{3}$ \& Mohamed Hanafi$^{1}$
\bigskip

{\it
$^{1}$ Oniris, INRAE, StatSC, Nantes, France. \\
$^{2}$ Epidemiology and Welfare, ANSES, Ploufragan, France. \\
$^{3}$ Department of Methodology and Statistics, Tilburg University, Tilburg, The Netherlands.
}
\end{center}
\bigskip

%--------------------------------------------------------------------------

\section{Introduction}

Data sets structured into blocks of variables measured on the same set of observations are frequently encountered in practice. These blocks may consist of quantitative variables originating from distinct measurement sources or correspond to thematic groupings defined by the user. For instance, in omics studies, blocks typically comprise variables derived from different analytical techniques (e.g., proteomics, metabolomics). In psychology, multiblock data structures commonly arise when multiple sets of variables are collected on the same individuals to capture different psychological domains. For example, distinct blocks may include self-report questionnaire scores assessing personality traits, behavioral task performance measures, and neurophysiological indicators such as electroencephalography  features.

The four expected objectives of exploratory multiblock data analysis are: (i) to perform a joint dimensionality reduction, (ii) to examine relationships among blocks, (iii) to investigate associations among variables both within and across blocks, and (iv) to capture within-block variability in order to reveal shared and block-specific structures in a common latent space.  An unified algebraic framework based on factorization strategies, denoted \textit{canonical factorisation}, was proposed by \cite{Bougeard024} for multiblock methods to achieve these objectives. In practice, canonical factorisation provides a structured framework for analysing data variability at multiple levels: differences between and within blocks, differences between observations by projecting them onto a common subspace. It also facilitates the harmonisation of outputs produced by different methods and supports a more objective criterion-based selection among competing approaches. 

This framework is a compass among multiblock models:  Generalized Canonical Correlation Analysis (GCCA, \cite{Carroll68}) and its popular case MAXVAR \cite{Carroll68}, Hierarchical Principal Component Analysis (HPCA, \cite{Wold96,Westerhuis98}), Structuration de Tableaux A Tois Indices de la Statistique (Statis, \cite{Lavit94}), Multiple Factorial Analysis (MFA, \cite{escofier1994}), Common Components and Specific Weights Analysis (CCWSA, \cite{Qannari001}), Multiple CO-inertia Analysis (MCOA, \cite{Chessel96})  fall within this framework. A benchmark of these methods had been done in \cite{Bougeard024}.

Initially, MCOA was proposed as an alternative to generalized canonical analysis, also known as MAXVAR \cite{Carroll68,Horst59}, for correlating multiple data tables. Compared with other generalisations of Hotelling's canonical correlation analysis \cite{Hotelling35}, such as those proposed by \cite{Kettenring71}, MCOA offers three main advantages. Firstly, it is computationally straightforward as it reduces to a singular value decomposition, the other approaches typically rely on iterative optimisation procedures. Secondly, MCOA is particularly effective at summarizing results across multiple datasets. Finally, MCOA provides sequential estimation of canonical factorisation under an orthogonality constraint.  In this work, we propose a new method, Optimal Multiple CO-inertia Analysis (OMCOA), which is a simultaneous version of MCOA. We compare this model with the previous one using simulated data whose structure is known.     

% \section{Canonical factorisation}

% Consider $N$ individuals characterized by $K$ blocks  $(\textbf{X}_1,\ldots,\textbf{X}_K)$ respectively of $p_1,\ldots,p_K$ variables. Let $\textbf{X}=[\textbf{X}_1|\ldots|\textbf{X}_K]\in \mathcal M_{N,P}(\mathbb R)$ ($P=\sum_{k=1}^K p_k$) be the concatenation of these $K$-blocks. Without loss of generality, the blocks are assumed to be column-centered.

% The $(K+1)-$blocks $(\textbf{X}_1,\ldots,\textbf{X}_K,\textbf{X})$ can be factorized in dimension $H=rank(\textbf{X})$ as follows  (see \cite{Bougeard024} for the demonstration) 
% \begin{align}
%   \left\{
% 	\begin{array}{ll}
% 		&\textbf{X}   = \textbf{T} \textbf{D} \textbf{P}^t   \\
% 		&\textbf{X}_k = \textbf{T} \textbf{D}_k \textbf{P}_k^t \quad \textrm{for} \quad k=(1, \dots, K)   \label{Eqn_Facto} \\	
% 		&\textrm{with the constraints} \quad \textbf{T}^t \textbf{T}=\textbf{I} \quad \textrm{and} \quad \|\textbf{p}^{(h)}\|=\|\textbf{p}_k^{(h)}\|=1 \quad \textrm{for} \quad h=(1, \dots, H)\ , 
% 	\end{array}
% 	\right.
% \end{align}
% where
% \begin{itemize}
% 	\item $\textbf{T}\in \mathcal M_{N,H}(\mathbb R)$ is the standardized and orthogonal common components,
% 	\item $\textbf{P} \in \mathcal M_{P,H}(\mathbb R)$ is the standardized (not necessarily orthogonal) common loadings,
% 	\item $\textbf{P}_k \in \mathcal M_{p_k,H}(\mathbb R)$ the standardized block-loadings, $\textbf{D}\in \mathcal M_{H,H}(\mathbb R)$
% 	\item $\textbf{D}_k\in \mathcal M_{H,H}(\mathbb R)$ the diagonal scaling matrices with property $\textbf{D}^2 = \sum_{k=1}^K \textbf{D}_k^2$ that relate common and block-analyses.
% \end{itemize}

% \noindent This factorization is \textit{the canonical factorization}. Each multiblock method derives a common component $\T$ using a different optimization criterion. Projecting the data onto these common components enables their canonical factorization (for details see \cite{Bougeard024}). The benchmark in \cite{Bougeard024} showed that all the methods under study produced similar results. The decomposition of the concatenated block $\X$ is unique for all the multiblock methods, however, only the MCOA method guarantees uniqueness of the decomposition of each block $\X_k$ (thanks to the orthoganlity constraint), and the proportion of variance it explains can be lower than that of the other methods. To ensure uniqueness of the solution and to improve the robustness of MCOA, we propose a simultaneous version.
 
%For $H<rank(\textbf{X})$, the objective is to determine the optimal approximation of the form given in (1) by minimizing an ordinary least square criterion. The models in section 1 (GCCA, SUMCOR, HPCA, Statis, CCWSA, MCOA) provide solution of this problem (see \cite{Bougeard024} for details on these methods). 

%The MCOA algorithm involves sequentially updating the columns $\Tt^{(1)},\ldots,\Tt^{(H)},\pp^{(1)},\ldots,\pp^{(H)}$ by minimizing the Frobenius norm of the residual. By construction, the MCOA estimations $\PP_1,...,\PP_K$ are orthogonal. 

\section{Optimal Multiple CO-Inertia Analysis}

\subsection*{Notations}

\begin{itemize}
	\item $K$ blocks  $(\textbf{X}_1,\ldots,\textbf{X}_K)$ respectively of $p_1,\ldots,p_K$ variables on $N$ individuals.
	\item  $\textbf{X}=[\textbf{X}_1|\ldots|\textbf{X}_K]\in \mathcal M_{N,P}(\mathbb R)$ ($P=\sum_{k=1}^K p_k$) be the concatenation of these $K$-blocks.
	\item $\textbf{U}^\intercal$ designates the transposed matrix of $\textbf{U}$. 
	\item $\Sigma_k=diag(\sigma_k^{(1)},...,\sigma_k^{(R)})$ are diagonal matrices 
	\item $\Vert . \Vert$ correspond to the Frobenius norm.
\end{itemize}
 Without loss of generality, the blocks are assumed to be column-centered.

\subsection{Minimization problem}
\cite{Bougeard024} proposed a unified framework for the factorization of blocks 
$(\X_1,\ldots,\X_K)$ (canonical factorization). As with the singular value 
decomposition of a single block $\X$, the key insight is to determine an optimal 
rank-$R$ approximation of $\X_1,\ldots,\X_K$ of the form
\begin{align}
  \left\{
    \begin{array}{ll}
        &\X  = \U \Sigma \V^\intercal +\textbf{E},   \\
        &\X_k = \U \Sigma_k \V_k^\intercal +\textbf{E}_k 
        \quad \text{for} \quad k=1, \dots, K   \label{Eqn_Facto2} \\	 
    \end{array}
  \right.
\end{align}
with the constraints $\U^\intercal\U=\textbf{I}_R$ and $\Vert\v_k^{(r)}\Vert=1$ for $r=1, \dots, R$.

\noindent Multiblock methods provide a solution to this canonical factorization 
under different approximation criteria (e.g., maximization of covariance or 
minimization of the Frobenius norm). Projecting the data onto $\U$ captures the 
structure of the observations (for details, see \cite{Bougeard024}). The benchmark 
study in \cite{Bougeard024} showed that all methods under consideration yielded 
similar results. The decomposition of the concatenated block $\X$ is unique across 
all multiblock methods; however, only MCOA guarantees the uniqueness of the 
decomposition of each individual block $\X_k$ (by virtue of the orthogonality of 
the columns of $\V_k$, for all $k=1,\ldots,K$). This comes  at the cost of explaining a lower proportion of variance than the other methods.

\subsection{Multiple CO-Inertia Analysis}
\cite{Chessel96} proposed Multiple Co-inertia Analysis as an alternative to MAXVAR \cite{Carroll68}. The rank-$H$ approximation is derived from successive rank-1 approximations of deflated versions of the blocks $\X_1,\ldots,\X_K$.

The rank-1 approximation criterion is based on minimising the sum of squared Frobenius norms
\begin{equation}
    \sum_{k=1}^K\Vert \X_k - \sigma_{k}\u\v_k^\intercal \Vert^2, \quad \Vert \u\Vert = \Vert \v_1\Vert = \ldots = \Vert \v_K\Vert = 1, \quad \sigma_k > 0. \label{mcoa_c}
\end{equation}

\noindent The solution to this minimisation problem is:
\begin{itemize}
  \item $\u$ is the first left singular vector of the concatenated block $\X$,
  \item $\v_k = \frac{\X_k^\intercal \u}{\Vert \X_k^\intercal \u \Vert}$,
  \item $\sigma_k = \Vert \X_k^\intercal \u \Vert$ for $k = 1, \ldots, K$.
\end{itemize}
 

The algorithm proposed by \cite{Chessel96} proceeds as follows: at each step $r = 1, \ldots, R$, the rank-1 approximation is applied to the deflated blocks $\X_k^{(r)} = \X_k(I_{p_k}-\sum_{i=1}^{r-1}\v_k^{(i)}{\v_k^{(i)}}^\intercal)$. Moreover, they proved that the resulting vectors $\u^{(1)}, \ldots, \u^{(R)}$ are mutually orthogonal. Note that, thanks to the deflation procedure,  $\v_k^{(1)}, \ldots, \v_k^{(R)}$ are mutually orthogonal for each $k$.
MCOA belongs to the class of so-called sequential multiblock component methods. In the first stage, block loadings are computed as solutions to an optimization criterion; higher-order block components are then computed from the deflated versions of the original blocks. The sequential approach is considered sub-optimal: 
although the parameter estimation at each stage is optimal in the least-squares sense, the corresponding criterion does not account for the full set of block components. The present study proposes an optimal alternative to MCOA.

\subsection{Optimal Multiple CO-inertia Analysis (OMCOA)}

Providing an optimal alternative to MCOA means minimizing
\begin{equation}
	\sum_{k=1}^K\Vert \X_k - \U\Sigma_k \V_k^\intercal\Vert^2= \Vert \X- \U \Sigma \V^\intercal\Vert^2, \label{Crit}
\end{equation}
with $\V_k^\intercal.\V_k=I_R$ for $k=1,\ldots,K$. Note that, with standardized and orthogonal block-loadings matrices $\V_1,\ldots,\V_K$, the solution of this problem is unique. 


\noindent Consider $\textbf{H}_k=\Sigma_k.\V^\intercal_k\in \mathcal M_{R,p_k}(\mathbb R)$ and $\textbf{H}=[\textbf{H}_1|\ldots|\textbf{H}_K]\in \mathcal M_{R,P}(\mathbb R)$. The iterative algorithm for building the solution of \eqref{Crit} requires to update successively $\U,\V_k,\Sigma_k$ as follows:
\begin{enumerate}
	\item \underline{Updating of $\U$:}  For given $\textbf{X}\in \mathcal M_{N,P}(\mathbb R)$, $\textbf{H}\in \mathcal M_{R,P}(\mathbb R)$, update $\U$ by $\U=\textbf{A}.\textbf{B}^\intercal$ where $\textbf{A},\textbf{B}$ are respectively the matrices of the first $R$ left singular vectors and the matrix of the first $R$ right singular vectors from the svd of $\textbf{X}.\textbf{H}^\intercal$.
	\item \underline{Updating of $\V_k$:} For given $\textbf{X}_k\in \mathcal M_{N,p_k}(\mathbb R)$, $\U\in \mathcal M_{R,N}(\mathbb R)$ with orthogonal columns, $\Sigma_k\in \mathcal M_{R,R}(\mathbb R)$, update $\V_k$ by: $\V_k=\textbf{A}\textbf{B}^\intercal$  where $\textbf{A},\textbf{B}$ are respectively the matrix of the first $R$ left singular vectors and the matrix of the first $R$ right singular vectors from the svd of $\textbf{X}^\intercal_k.\U$.
	 \item \underline{Updating of $\Sigma_k$:}  For given $\textbf{X}_k\in \mathcal M_{N,p_k}(\mathbb R)$, $\U\in \mathcal M_{R,N}(\mathbb R)$ orthogonal, $\V_k\in \mathcal M_{R,p_k}(\mathbb R)$ with orthonormal columns, $\Sigma_k$ is given by 
	 \[\Sigma_k=diag(\U^\intercal\textbf{X}_k\V_k).\] 
\end{enumerate}

\begin{proposition}
The successive updates of $\U$, $\V_k$, and $\Sigma_k$ guarantee that the criterion defined in \eqref{Crit} decreases or remains constant at each step, ensuring convergence of the algorithm.
\end{proposition}


% \begin{lemma} \label{lemma1}
% 	For given matrices $\X\in\mathcal M_{N,P}(\mathbb R),\h\in\mathcal M_{P,R}(\mathbb R)$, the maximum of $\trace(\T^\intercal\X\h)$ with the constraint $\T^\intercal.\T=I_R$ is $\T=\textbf{U}.\textbf{V}^\intercal$ where the columns of $\textbf{U}$ (respectively $\textbf{V}$) are the $R$ first left (respectively right) singular vectors of $\X\h$. (see \cite{cliff66}) 
% \end{lemma}

% \begin{proof}
% 	Consider the SVD of $\X.\h =\textbf{U}.\Sigma.\textbf{V}^t$. We have $\trace(\textbf{T}^t.\textbf{U}.\Sigma.\textbf{V}^t)=\trace(\Sigma.\textbf{V}^t.\textbf{T}^t.\textbf{U})$
% 	($\trace(\textbf{A}.\textbf{B})=\trace(\textbf{B}.\textbf{A})$). The matrix $\textbf{W}=\textbf{V}^t.\textbf{T}^t.\textbf{U}$ is orthogonal. 
% The problem of maximization remains to maximizing 
% \[\trace(\Sigma.\textbf{W})=\sum_{h=1}^H\sigma^2_h.w_{hh}.\]
% Since $\textbf{W}$ is orthogonal, $|w_{hh}| \leq 1$ for all $h$. The trace of $\Sigma.\textbf{W}$ is maximal for $w_{hh}=1$, that implies $\textbf{W}=I_H$ and $\textbf{T}=\textbf{U}.\textbf{V}^t$. Moreover, the maximum of the trace is  $\sum_{h=1}^H\sigma^2_h$. 

% \end{proof}

% \begin{lemma} \label{lemma2}
% 	For given matrices $\X\in\mathcal M_{N,P}(\mathbb R),\h\in\mathcal M_{P,R}(\mathbb R), \T \in \mathcal M_{N,R}(\mathbb R)$ we have 
% 	\[\trace(\textbf{X}^\intercal.\textbf{T}.\textbf{H}^\intercal)=\trace(\textbf{T}^\intercal\textbf{X}\textbf{H}).\]
% \end{lemma}

%  \begin{proof}
%  $\trace(\X^\intercal\T\h^\intercal)=\trace(\h^\intercal\X^\intercal\T)=\trace(\T^\intercal\X\h).$
%  The first equality comes from $\trace(\textbf{A}\textbf{B})=\trace(\textbf{A}\textbf{B})$ for $\textbf{A},\textbf{B}$ of dimension $N,P$ and $P,N$. The second equality comes from the equality of the trace of a square matrix and its transpose.
%  \end{proof}


\medskip
\noindent
\textbf{Remark.} This proof guarantees convergence of the values of the criterion, but not necessarily uniqueness of the limit point. Uniqueness of the solution is established separately, via the orthogonality of the columns of $\V_k$ for all $k$.

\section{Comparison of MCOA and OMCOA}
\subsection{Simulation study}
MCOA and OMCOA  are compared on simulated data. Three blocks $\X_1,\X_2,\X_3$ with respectively $p_1=10,p_2=15,p_3=20$ variables and $I=100$ observations have been generated from:
\[\X_k=\U\Sigma_k\V_k^T+\textbf{E}_k, \quad \text{ for } k=1,\ldots,3,\]  
where $\textbf{E}_k$ is the residual term and for $k=1,\ldots,3$ the columns of the matrices $\U_k$ (respectively $\V_k$) are orthonormal.
%Note that, when the blocks are centered by columns, the square of diagonal terms $d_{rk}=d_r(\X_k)$ corresponds to the variance of the block $\X_k$ explained by the common dimension $\textbf{t}_r$. 
We compare the estimation of $\Sigma_1,...\Sigma_K$ obtained using MCOA and OMCOA.
We consider the following simulation scenario

$$\Sigma_k(t)=(1-t)\,\Sigma^{(0)}_k+t\,\Sigma_k^{(1)}, \quad t\in[0,1].$$ 
where :
\begin{itemize}
    \item For $t=0$, all the blocks have a common profile. \\
    $\Sigma^{(0)}_1=\Sigma^{(0)}_2=\Sigma^{(0)}_3=\begin{pmatrix}
        0.7 & & \\
        & 0.5 & \\
        & & 0.3
    \end{pmatrix}$
   
    \item  For $t=1$ each block has its own dominant dimension of the common space $\U$ with permuted profiles.\\
     $\Sigma^{(1)}_1=\begin{pmatrix}
        0.7 & & \\
        & 0.5 & \\
        & & 0.3
    \end{pmatrix}$, $\Sigma^{(1)}_2=\begin{pmatrix}
        0.5 & & \\
        & 0.3 & \\
        & & 0.7
    \end{pmatrix}$, $\Sigma^{(1)}_3=\begin{pmatrix}
        0.3 & & \\
        & 0.7 & \\
        & & 0.5
    \end{pmatrix}$.
\item  $\sum_{k,r}(\sigma_k^{(r)}(0))^2=\sum_{k,r}(\sigma_k^{(r)}(1))^2=0.83$.
\item In all cases, $\Vert \EE_k\Vert^2=.10 \times \Vert \X_k\Vert^2$
\end{itemize} 
 


\subsection{Results}

Figure \ref{fig_s1} and  \ref{fig_s2} provide the results of the estimation by MCO and OMCOA respectively for the two scenari described in Section 5.1.
As expected, the estimates in the common-configuration case were very close (see Figure \ref{fig_s1}). In the second scenario, where all the blocks had different configurations, OMCOA outperformed MCOA in terms of both bias and precision across all parameter sets (see Figure \ref{fig_s2}). In both scenari, the estimates obtained with OMCOA did not vary with the magnitude of the residuals.

% \section*{Conclusion}
% \cite{Bougeard024} have compared estimations for the canonical factorisation problem. In this work, we propose on optimal version of Multiple COinertia Analysis that performs better on the simulated data. Constraining the matrices $(\PP_1,\dots,\PP_K)$ to be orthogonal ensures uniqueness of the solution. In particular, if $|\mathbf{E}_k|=0$, the OMCOA method recovers the exact canonical factorization while other methods do not guarantee this."    
% Furthermore, in practical applications the canonical factorisation is more readily interpreted under orthogonal configurations, as these yield a clearer decomposition of the effects.

\begin{figure}[hbt]
  \includegraphics[width=\linewidth]{plot1.png}
  
\end{figure}


\begin{figure}[hbt]
  \includegraphics[width=\linewidth]{plot2.png}
\end{figure}


\bibliographystyle{apalike}
\bibliography{Biblio}

% \begin{appendix}


% \begin{proof}
% 	The criterion to be minimized is:
% \[
%   \mathcal{L}(\T, \PP_1, \ldots, \PP_K, \D_1, \ldots, \D_K)
%   = \sum_{k=1}^{K} \bigl\| \X_k - \T \D_k \PP_k^\top \bigr\|^2,
% \]
% subject to $\T^\intercal \T = I_H$ and $\PP_k^\top \PP_k = I_H$ for all $k = 1, \ldots, K$.

% The algorithm is of the \emph{Alternating Least Squares} (ALS) type: at each iteration, the criterion is minimized successively with respect to each parameter, while the others are held fixed. It therefore suffices to show that each individual update decreases or maintains $\mathcal{L}$.

% \paragraph{Update of $\T$} With $\PP_k$ and $\D_k$ fixed, let $\h_k = \D_k \PP_k^\top \in \mathcal{M}_{R,p_k}(\mathbb{R})$ and
% $\h = [\h_1 \mid \cdots \mid \h_K] \in \mathcal{M}_{R,P}(\mathbb{R})$. Then:
% \[
%   \mathcal{L} = \sum_{k=1}^{K} \|\X_k - \T \h_k\|^2 = \|\X - \T \h^\intercal\|^2.
% \]
% Expanding:
% $\|\X - \T \h^\intercal\|^2
%   = \|\X\|^2 - 2\,\operatorname{trace}(\X^\intercal \T \h^\intercal) + \|\T \h^\intercal\|^2.$

% \noindent Since $\T^\intercal \T = I_H$: $\|\T \h^\intercal\|^2
%   = \operatorname{trace}(\h \T^\intercal T H^\intercal)
%   = \operatorname{trace}(\h \h^\intercal)
%   = \|\h\|^2 \quad \text{(constant w.r.t.\ } \T\text{)}.$

% \noindent Minimizing $\mathcal{L}$ with respect to $\T$ is therefore equivalent to
% maximizing $$\trace(\X^\intercal\T\h^\intercal)=\trace(\h^\intercal\X^\intercal\T)=\trace(\T^\intercal\X\h),$$
% subject to $\T^\intercal \T = I_R$. The first equality comes from $\trace(\textbf{A}\textbf{B})=\trace(\textbf{A}\textbf{B})$ for $\textbf{A},\textbf{B}$ of dimension $N,P$ and $P,N$. The second equality comes from the equality of the trace of a square matrix and its transpose. By \cite{cliff66}, the maximum is attained at $\T^* = UV^\intercal$, where $U$ (resp.\ $V$) are the matrices of the $R$ leading left (resp.\ right) singular vectors of $\X\h$.
% This update is the global minimizer of $\mathcal{L}$ with respect to $\T$, hence:
% \[
%   \mathcal{L}(\T^*, \ldots) \leq \mathcal{L}(\T_{\mathrm{old}}, \ldots).
% \]
% \paragraph{Update of $\PP_k$} With $\T$ and $\D_k$ fixed, expanding the contribution of block $k$:
% \[
%   \|\X_k - \T \D_k \PP_k^\intercal\|^2
%   = \|\X_k\|^2
%     - 2\,\operatorname{trace}(\X_k^\intercal \T \D_k \PP_k^\intercal)
%     + \|\T \D_k \PP_k^\intercal\|^2.
% \]
% Since $\T^\intercal \T = I_R$ and $\PP_k^\intercal \PP_k = I_R$:
% \[
%   \|\T \D_k \PP_k^\intercal\|^2
%   = \operatorname{trace}(\PP_k \D_k^\intercal \T^\intercal \T \D_k \PP_k^\intercal)
%   = \operatorname{trace}(\D_k^2)
%   \quad \text{(constant w.r.t.\ } \PP_k\text{)}.
% \]
% Minimizing therefore reduces to maximizing
% $\operatorname{trace}(\X_k^\intercal \T \D_k \PP_k^\intercal) = \operatorname{trace}(\PP_k^\intercal\X_k^\intercal\T\D_k)$
% subject to $\PP_k^\intercal \PP_k = I_R$. Since the diagonal matrix $\D_k > 0$ does not affect the singular vectors, the maximum is attained at $\PP_k^* = \textbf{UV}^\intercal$ where $\textbf{U}$, $\textbf{V}$ are obtained from the SVD of $\X_k^\intercal\T$ \cite{cliff66} . Hence:
% \[
%   \mathcal{L}(\ldots, \PP_k^*, \ldots) \leq \mathcal{L}(\ldots, \PP_{k,\mathrm{old}}, \ldots).
% \]

% %---------------------------------------------------------------
% \paragraph{Update of $\D_k$} With $\T$ and $\PP_k$ fixed (both having orthonormal columns), and $\D_k = \operatorname{diag}(d_1,
% \ldots, d_R)$, the contribution of block $k$ is:
% \[
%   \|\X_k - \T \D_k \PP_k^\intercal\|^2
%   = \|\X_k\|^2
%     - 2\,\operatorname{trace}(\PP_k^\intercal \X_k^\intercal \T \D_k)
%     + \operatorname{trace}(\D_k^2).
% \]
% Since $\D_k$ is diagonal, it remains to minimazing:
% \[
%   \sum_{r=1}^{R} \Bigl[ d_r^2 - 2\,(\T^\intercal \X_k \PP_k)_{rr} \cdot d_r \Bigr].
% \]
% Each term is a convex parabola in $d_r$, minimized at: $d_r^* = (\T^\intercal \X_k \PP_k)_{rr},$ or, in matrix form:$\D_k^* = \operatorname{diag}(\T^\intercal \X_k \PP_k).$
% This update is therefore the exact minimizer with respect to $\D_k$, and:
% \[
%   \mathcal{L}(\ldots, \D_k^*, \ldots) \leq \mathcal{L}(\ldots, \D_{k,\mathrm{old}}, \ldots).
% \]

% %---------------------------------------------------------------
% \paragraph{Global Convergence} At each complete iteration (update of $\T$, then all $\PP_k$, then all $\D_k$), the criterion $\mathcal{L}$ is monotonically decreasing and bounded below by zero. By the monotone convergence theorem, the sequence of values $(\mathcal{L}^{(t)})_{t \geq 0}$
% converges:
% \[
%   \mathcal{L}^{(t+1)} \leq \mathcal{L}^{(t)}, \quad \mathcal{L}^{(t)} \geq 0
%   \implies
%   \mathcal{L}^{(t)} \xrightarrow[t \to \infty]{} \ell^* \geq 0.
% \]
% \end{proof}

% \end{appendix}


%\printbibliography

\end{document}

